\chapter{Implementation}
\label{chap:implementation}

This chapter describes the implementation of all classifiers evaluated in this work. The full source code is publicly available~\cite{motamedifar2026repo} and organised in the \texttt{src/} directory, written in Python~3.11.

% ============================================================
\section{Software Environment}
\label{sec:env}

Three main libraries are used. \textbf{PennyLane}~\cite{pennylane} (v0.39) provides the QNode interface for circuit evaluation, the \texttt{lightning.qubit} noiseless simulator used in the main experiments, and the \texttt{AdamOptimizer}; gradients are computed with adjoint differentiation for speed. \textbf{scikit-learn}~\cite{sklearn} provides the classical classifiers and model-selection utilities. \textbf{pandas} and \textbf{NumPy} handle data loading and numerical operations. All experiments run on CPU under noiseless simulation; no real quantum hardware is used.

% ============================================================
\section{Classical Baselines}
\label{sec:impl_classical}

The classical baselines are implemented in \texttt{src/classical/classical\_baselines.py} and evaluated under both experimental settings defined in Section~\ref{sec:splits}.

\subsection{Support Vector Machine}

The SVM uses an RBF kernel. Hyperparameters $C \in \{0.1,\,1,\,10,\,100\}$ and $\gamma \in \{\texttt{scale},\,\texttt{auto},\,0.01,\,0.1\}$ are selected by 5-fold cross-validation on the training set. The best configuration is $C = 100$, $\gamma = \texttt{scale}$, achieving a cross-validation accuracy of 97.8\%.

\subsection{Random Forest}

A Random Forest of 100 decision trees is trained with default scikit-learn settings. As an ensemble method, it provides a robust baseline for tabular data and additionally yields feature importance scores that complement the ANOVA-based selection in the preprocessing stage.

\subsection{Neural Networks}

Two fully connected networks with ReLU activations are trained using Adam ($\eta = 10^{-3}$, up to 200 epochs):
\begin{itemize}[nosep]
    \item \textbf{NN-Small} ($8 \to 12 \to 6 \to 3$, ${\sim}207$ parameters): a small network in the same order of magnitude as the quantum classifiers.
    \item \textbf{NN-Medium} ($8 \to 16 \to 8 \to 3$, ${\sim}307$ parameters): a slightly larger network to test whether additional classical capacity helps.
\end{itemize}

\subsection{Autoencoder (Unsupervised Anomaly Detector)}
\label{sec:impl_autoencoder}

Supervised classifiers learn the boundary between the classes they are shown during training, which is precisely why they fail on unseen attack types. As a fairer classical baseline for the zero-day setting I add a shallow \emph{autoencoder} trained on normal traffic only. The architecture is a symmetric $8 \to 6 \to 3 \to 6 \to 8$ MLP with $\tanh$ activations and Xavier initialization, optimized with Adam ($\eta = 2 \times 10^{-3}$, weight decay $10^{-5}$, cosine learning-rate schedule) over 100~epochs on the ${\sim}28{,}000$ normal training samples. At inference, an input is flagged as anomalous if its reconstruction MSE exceeds the 95\textsuperscript{th} percentile of the training reconstruction errors. The threshold is set to the 95\textsuperscript{th} percentile of the training-normal reconstruction errors; the held-out normal split is used only to estimate the resulting false-positive rate, which is approximately 5\%. This protocol turns ``zero-day detection'' into ``anything that does not look like normal traffic,'' and is the standard unsupervised approach to network intrusion detection.

Table~\ref{tab:model_params} gives an overview of all models and their parameter budgets.

\begin{table}[ht]
\centering
\begin{tabular}{llc}
    \toprule
    \textbf{Model} & \textbf{Type} & \textbf{Parameters} \\
    \midrule
    SVM (RBF, $C{=}100$)          & Classical & kernel-based \\
    Random Forest (100 trees)      & Classical & tree-based   \\
    NN-Small $(12,6)$              & Classical & $\sim$207    \\
    NN-Medium $(16,8)$             & Classical & $\sim$307    \\
    Autoencoder $(8\!-\!6\!-\!3\!-\!6\!-\!8)$ & Classical & $\sim$160    \\
    \midrule
    Data Re-uploading (1 qubit, 6 layers)  & Quantum & 108 \\
    Data Re-uploading (4 qubits, 3 layers) & Quantum & 216 \\
    QCNN (8 qubits, 3 stages)             & Quantum & 81  \\
    \bottomrule
\end{tabular}
\caption{Parameter counts for all evaluated models. Quantum counts are the total for the 3-class one-vs-all system ($3\times$ the per-classifier count).}
\label{tab:model_params}
\end{table}

% ============================================================
\section{Quantum Classifiers}
\label{sec:impl_quantum}

Both quantum classifiers are implemented in \texttt{src/quantum/} using PennyLane's \texttt{lightning.qubit} simulator with adjoint differentiation, which gives identical results to \texttt{default.qubit} but runs noticeably faster. They share the same one-vs-all strategy: for $K = 3$ classes, three independent binary circuits are trained, each targeting one class with labels $y_i \in \{-1, +1\}$. At inference, the class whose circuit returns the highest $\langle Z \rangle$ is predicted.

\subsection{QCNN}

The QCNN operates on 8~qubits with 3~conv-pool stages ($8 \to 4 \to 2 \to 1$), followed by an $R_z R_y R_z$ fully connected layer on the remaining qubit. The convolutional unitary is a 6-parameter two-qubit block inspired by the ansatz family of Hur et al.~\cite{hur2022quantum}:
\[
U(\boldsymbol{\theta}) = \big[R_y R_z\big]_{q_1} \;\text{CNOT}\; \big[R_y R_z\big]_{q_0} \big[R_y R_z\big]_{q_1},
\]
a lighter variant of Hur et al.'s 10-parameter $U_6$ (which uses $R_x/R_z$ rotations and CRX entanglement). Each pooling layer uses a CRZ--CRY pair (2~parameters) to conditionally transfer information before discarding one qubit, a simplification of the CRZ + controlled-$R_x$ pooling proposed in the original paper. The total per-classifier parameter count is $3 \times 6 + 3 \times 2 + 3 = 27$, and $81$~parameters for the full 3-class system.

Input data is encoded via angle encoding: each of the 8~features is mapped to an $R_y$ rotation on the corresponding qubit.

\subsection{Data Re-uploading}

Two variants are implemented:
\begin{itemize}[nosep]
    \item \textbf{Single-qubit}: 6~layers on 1~qubit, $6 \times 6 = 36$~parameters per classifier (108~total).
    \item \textbf{Multi-qubit}: 3~layers on 4~qubits with nearest-neighbor CZ entanglement, $6 \times 4 \times 3 = 72$~parameters per classifier (216~total).
\end{itemize}

In each layer, rotation angles are computed as $\theta_j^{(l)} + w_j^{(l)} x_{f(j,l)}$, where the feature index $f(j,l)$ cycles through the input features across layers. For the multi-qubit variant, the output is the average $\langle Z \rangle$ over all qubits, expressed as a single Hamiltonian measurement $\langle H \rangle$ with $H = \frac{1}{n}\sum_q Z_q$.

\subsection{Training Configuration}

All quantum classifiers are trained with the Adam optimizer (learning rate $\eta = 0.01$) for 100~epochs. Each one-vs-all binary classifier uses a \emph{separate} optimizer instance to avoid momentum interference between sub-models. Gradients are computed via adjoint differentiation, which is mathematically equivalent to the parameter-shift rule but considerably faster in simulation. Parameters are initialized uniformly in $[-0.1, 0.1]$ with a controlled random seed to ensure reproducibility and to mitigate barren plateaus~\cite{mcclean2018barren}. To account for the known sensitivity of quantum circuits to initialization, the main experiments are repeated with 10~different seeds and results are reported as mean~$\pm$~standard deviation (Section~\ref{sec:res_sensitivity}). For the standard 3-class task, training uses a stratified subsample of 500~samples and test evaluation uses 200~stratified samples; for zero-day evaluation, the same sizes are drawn as fixed (non-stratified) random subsamples of the normal+DoS train set and the injection-only test set. All circuits run on noiseless simulation.

% ============================================================
\section{Evaluation Protocol}
\label{sec:eval_protocol}

All models are evaluated under the two settings defined in Section~\ref{sec:splits}:
\begin{itemize}[nosep]
    \item \textbf{Standard}: 3-class accuracy and weighted F1-score on the 200-sample test set.
    \item \textbf{Zero-day}: binary detection rate on the injection-only test set (7,034~samples after the 80/20 split), where the model has only seen normal and DoS traffic during training. Quantum models are evaluated on a fixed 200-sample subsample of this set due to simulation cost.
\end{itemize}
Results are reported in Chapter~\ref{chap:results}.
